\chapter*{Resumo}

% =====================================================================
% Traducao fiel do Abstract (Cap00/Abstract.tex). Escrever DEPOIS do
% Abstract, para nao divergirem. Mantido em portugues porque programas de
% pos-graduacao brasileiros geralmente exigem o resumo no vernaculo mesmo
% quando o trabalho e escrito em ingles -- confirmar a exigencia do PPGI
% com o orientador. Se nao for exigido, apagar este arquivo e a chamada
% \resumo em main.tex.
% =====================================================================

\noindent Modelos de difusão geram faces fotorrealistas a partir de texto, mas o controle sobre a
\emph{intensidade} de uma expressão ainda é grosseiro: comandos textuais são categóricos,
distinguindo ``sorrindo'' de ``não sorrindo'' em vez de comandar uma intensidade graduada. Os
sistemas mais próximos deste problema condicionam a geração a intensidades contínuas de unidades
de ação, mas produzem cada nível de intensidade como uma imagem independente, não impõem nem medem
a identidade entre os níveis, e relatam o controle de intensidade apenas qualitativamente. Este
trabalho reformula a intensidade da expressão facial como um parâmetro de geração em nível de
sequência: dado um texto e uma trajetória de intensidade por quadro, um \emph{backbone} Stable
Diffusion~1.5 congelado, estendido com um módulo de movimento AnimateDiff, gera um único clipe
temporalmente coerente cuja intensidade de expressão por quadro segue o comando. Apenas um pequeno
codificador de expressão e os pesos de atenção que integram sua saída à atenção temporal do
\emph{backbone} são treinados; o \emph{backbone} permanece congelado. Resultados preliminares em um
único eixo de sorriso, a unidade de ação AU12 (levantadora do canto do lábio), mostram que a
intensidade detectada segue a trajetória comandada (correlação de Pearson $r = 0{,}77$, confirmada
por um detector independente que não produziu nenhum dos rótulos de treinamento) e que a identidade
é preservada ao longo do clipe melhor do que pela geração independente (similaridade facial mínima
entre quadros adjacentes de $0{,}88$ contra $0{,}79$). A comparação direta com o trabalho anterior
mais próximo empata no seguimento da trajetória. Relatado de forma honesta, o controle é relativo à
trajetória comandada, e não absoluto, consequência de dados de treinamento do tipo início-ao-ápice
que o trabalho anterior oculta por não medir. A contribuição é, portanto, a formulação nativa em
sequência da intensidade da expressão, juntamente com um protocolo de avaliação que separa
seguimento de trajetória, calibração absoluta e identidade entre quadros, e que trata a
circularidade rótulo--avaliação que os trabalhos mais próximos deixam sem sinalizar.

\vspace{1em}
\noindent\textbf{\textit{Palavras-chave}:~} síntese de expressões faciais. modelos de difusão. unidades de ação. geração controlável. consistência temporal.
